\documentclass{article}
\usepackage[UTF8]{ctex}
\usepackage{amsmath,amssymb,amsfonts}
\usepackage{enumitem}
\usepackage{geometry}
\geometry{a4paper, margin=1in}
% === 错题本排版优化（不参与题目标识） ===
\setlength{\parindent}{0pt}
\setlength{\parskip}{0.35em plus 0.10em minus 0.08em}
\setlength{\emergencystretch}{3em}
\setlist{itemsep=0.25em, topsep=0.40em, parsep=0pt}
\allowdisplaybreaks[2]
% === 错题本排版优化结束 ===
\begin{document}
\section*{东北大学2024年数学分析真题}
\section*{72．东北大学2024年数学分析试卷}

\subsection{题目1}
1．（15 分）求极限 \(\lim _{x \rightarrow 0} \frac{\frac{x}{1+x}-\ln (1+x)}{x^{2}}\) ．

\subsection{题目2}
2．（15 分）已知 \(f(x)=\int_{1}^{x} e^{-y^{2}} \mathrm{~d} y\) ，求 \(\int_{0}^{1} x^{2} f(x) \mathrm{d} x\) ．

\subsection{题目3}
3．（15 分）求级数 \(\sum_{n=0}^{\infty} \frac{(-1)^{n}\left(n^{2}-n+1\right)}{2^{n}}\) 的和。

\subsection{题目4}
4．（15 分）求曲面积分

\[
I=\iint_{\Sigma}(x y+y z+z x) \mathrm{d} S
\]

其中 \(\sum\) 为锥面 \(z=\sqrt{x^{2}+y^{2}}\) 被柱面 \(x^{2}+y^{2}=2 a x\) 所截的部分。

\subsection{题目5}
5．（15 分）设 \(0 \leq \alpha \leq 1\) ，证明：\(f(x)=x^{\alpha}\) 在 \([0,+\infty)\) 上一致连续．

\subsection{题目6}
6．（15 分）证明：\(\sum_{n=1}^{\infty}\left(1-\cos \frac{x}{n}\right)\) 在 \((-a, a)\) 上一致收敛。

\subsection{题目7}
7．（15 分）设数列 \(\left\{x_{n}\right\}\) 有界且 \(\lim _{n \rightarrow \infty}\left(x_{2 n}+2 x_{n}\right)\) 存在．证明： \(\lim _{n \rightarrow \infty} x_{n}\) 也存在．

\subsection{题目8}
8．（15 分）证明：数列 \(x_{n}=\sum_{i=1}^{n}(-1)^{i-1} \frac{1}{i}\) 收敛。

\subsection{题目9}
9．（15 分）设 \(f(x) \in C[a, b]\) 且单调增加，证明不等式： \(\int_{a}^{b} x f(x) \mathrm{d} x \geq \frac{a+b}{2} \int_{a}^{b} f(x) \mathrm{d} x\) ．

\subsection{题目10}
10．（15 分）设 \(f(x)\) 是定义在有限闭区间 \([a, b]\) 上的函数，且对任意的 \(x \in[a, b]\) ，有 \(a \leq f(x) \leq b\) ，同时存在 \(L \in(0,1)\) ，使得对任意的 \(x, y \in[a, b]\) ，有

\[
|f(x)-f(y)| \leq L|x-y| .
\]

对任意的 \(x_{0} \in[a, b]\) ，定义 \(x_{n}=f\left(x_{n-1}\right), n=1,2, \cdots\) 。证明： \(\lim _{n \rightarrow \infty} x_{n}=\xi \in[a, b]\) ，并且 \(f(\xi)=\xi\) 。
\end{document}
